% Provenance (never printed): see sources/outline_sources.md (Year 3:
% partition coefficients; batch and fractional melting; crystallisation
% models; reading a spider diagram quantitatively). Brief:
% sources/book5/BRIEFS.md, ch. 12. Curves from figdata/bachelor-3/
% melting_models.py. Exercise and problem values in the Book 5 scratch
% calculator calc/ch12.py. Aldren Island and its lavas are invented.
\chapter{微量元素模拟}\label{ch:b3:trace-element-modelling}

同一岛上两块相距数千米的玄武岩喷发，其中一块
的镧含量是另一块的一倍半以上，而镱含量则少了近四分之三。
它们是来自两个不同的地幔，还是来自同一地幔、
熔融量不同，还是熔融深度不同？主量元素
几乎无法说明问题：两者都是玄武岩。微量元素随矿物迁移
而不改变矿物，它们记录了生成了多少熔体、熔体如何离开
源区以及哪些矿物留在了残余中。本章将这一记录写成方程，
并由此读出其中信息。

\begin{recall}
矿物–熔体分配系数 $D$、总分配系数、
相容元素与不相容元素、稀土元素、CI
球粒陨石与原始地幔、标准化图解与铕
异常、MORB 与 OIB 见本科二年级分册中关于
地球中元素的一章。橄榄岩、二辉橄榄岩、部分熔融与熔体分数 $F$，
见其中关于地幔熔融的一章；平衡结晶与分离
结晶、瑞利分馏定律、混合与
同化–分离结晶（AFC），见其中关于岩浆
分异的一章。那里已证明，两种极不相容的
元素之比在结晶过程中保持不变。
\end{recall}

\section{熔融岩石的物质衡算}\label{sec:b3:trace-element-modelling:balance}

质量为 1 的一块岩石发生熔融。熔体分数为 $F$ 时，它生成了质量为 $F$ 的
熔体，并保留了质量为 $1 - F$ 的固体残余。某微量元素在
源区中的浓度为 $C_0$，它在二者间分配：熔体中的浓度为
$C_L$，残余中的浓度为 $C_S$。本章的每一个模型都是
同一平衡式的变形，$C_0 = FC_L + (1 - F)C_S$ 适用于整个熔体与残余，
并附加一条联系 $C_S$ 与 $C_L$ 的规则：在晶体与熔体的接触处，
$C_{\text{矿物}} = D\,C_L$。各模型的区别在于哪些相彼此接触，
以及接触发生在何时。

地幔中各矿物的分配系数相差很大。据
一份汇编，镧在单斜辉石中 $D = 0.058$，在
石榴子石、橄榄石、斜方辉石和尖晶石中均为 $0.001$；镱在单斜辉石中 $D = 0.58$，
在石榴子石中为 $4.5$，在其余三种矿物中为 $0.16$。石榴子石
容纳重稀土，排斥轻稀土。
% ledger: budahn91:d-ree

\begin{definition}[非实比熔融]\label{def:b3:trace-element-modelling:nonmodal}
在\emph{非实比熔融}\index{非实比熔融}中，源区各矿物
进入熔体的比例 $p_i$ 不同于它们在源区中的质量
分数 $x_i$。熔体起初所见的是源区的
总分配系数 $D_0 = \sum_i x_iD_i$；随熔融进行，它所见的则是
按熔融量加权的系数 $P = \sum_i p_iD_i$。
\end{definition}

熔融很少是实比的。在一个尖晶石二辉橄榄岩的算例中，
其组成为 \qty{58}{\percent} 橄榄石、\qty{27}{\percent} 斜方辉石、
\qty{12}{\percent} 单斜辉石和 \qty{3}{\percent} 尖晶石，各矿物
进入熔体的比例分别为 20、25、45 和 \qty{10}{\percent}：
单斜辉石熔化的速度是其份额的四倍，橄榄石则
慢至三分之一。
% ledger: white13:ex73

\begin{proposition}[残余的总分配系数]
\label{prop:b3:trace-element-modelling:residue}
在 $p_i$ 为常数的\omterm{def:b3:trace-element-modelling:nonmodal}{非实比熔融}中，熔体分数为 $F$ 时，残余的
总分配系数为
\[
\bar D(F) = \frac{D_0 - PF}{1 - F},
\]
上式在每种矿物均未耗尽时成立，即 $F < x_i/p_i$ 对每个 $i$ 均成立。
\end{proposition}

\begin{proof}
以单位质量源区计，矿物 $i$ 在 $F$ 时的剩余质量为 $x_i - p_iF$，
残余的质量为 $1 - F$。其总分配系数为按质量加权的
平均值，$\sum_i (x_i - p_iF)D_i/(1 - F) = (D_0 - PF)/(1 - F)$。当
$x_i - p_iF = 0$ 时，矿物耗尽。
\end{proof}

取上述矿物分配系数，对前述算例中的尖晶石二辉橄榄岩，
以及在实验室中于 \qty{5}{GPa} 和
\qty{1400}{\celsius} 下保持的一块石榴子石二辉橄榄岩（按质量计，橄榄石 \qty{54.1}{\percent}、
斜方辉石 \qty{9.8}{\percent}、单斜辉石 \qty{19.5}{\percent} 和
石榴子石 \qty{16.7}{\percent}），其总分配系数列于下表。
% ledger: white13:ex73, rfr25:gkr-5gpa, budahn91:d-ree
\begin{center}
\small
\begin{tabular}{lccc}
\hline
源区 & $D_{\mathrm{La}}$ & $D_{\mathrm{Sm}}$ & $D_{\mathrm{Yb}}$ \\
\hline
尖晶石二辉橄榄岩 & 0.0078 & 0.062 & 0.21 \\
石榴子石二辉橄榄岩 & 0.012 & 0.107 & 0.97 \\
\hline
\end{tabular}
\end{center}
镧在两种源区中都是极不相容的。镱在尖晶石二辉橄榄岩中为中等不相容，
在石榴子石二辉橄榄岩中则近于中性。
两种源区的差别几乎全部体现在重
稀土上。

\section{批式熔融}\label{sec:b3:trace-element-modelling:batch}

\begin{theorem}[批式熔融]\label{thm:b3:trace-element-modelling:batch}
在\emph{批式熔融}\index{批式熔融}中，熔体与
全部残余保持平衡，直至整批熔体一次性离开。当总分配系数
$D$ 为常数时，熔体分数为 $F$ 时，熔体与残余
中的浓度满足
\[
\frac{C_L}{C_0} = \frac{1}{D + F(1 - D)}, \qquad
\frac{C_S}{C_0} = \frac{D}{D + F(1 - D)}.
\]
\end{theorem}

\begin{proof}
平衡时，残余体中每个晶体含有 $D_iC_L$，
因此残余体平均含有 $C_S = DC_L$。质量平衡 $C_0 = FC_L + (1 - F)DC_L$ 可解得
$C_L$，进而得到 $C_S$。
\end{proof}

该式即本科二年级分册中平衡结晶定律倒过来读的结果：
两者中 $F$ 都是封闭系统中液体的比例。其两个
极限值得了解。当 $F \to 0$ 时，$C_L \to C_0/D$：第一批熔体
至多使某元素富集 $1/D$ 倍。当 $D \to 0$ 时，$C_L \to
C_0/F$：完全不相容元素被稀释进熔体，其
富集程度仅取决于 $F$。对于相容元素，$D > 1$，熔体
贫于源区，且 $F$ 越小，贫化越甚。

\begin{proposition}[非实比批式熔融]\label{prop:b3:trace-element-modelling:nonmodal}
在\omterm{thm:b3:trace-element-modelling:batch}{批式熔融}中，熔融比例 $p_i$ 按\cref{def:b3:trace-element-modelling:nonmodal}定义，
\[
\frac{C_L}{C_0} = \frac{1}{D_0 + F(1 - P)}.
\]
\end{proposition}

\begin{proof}
残余体含有 $\bar D(F)\,C_L$（见\cref{prop:b3:trace-element-modelling:residue}），
故 $C_0 = FC_L + (1 - F)\bar D(F)C_L = C_L(F + D_0 - PF)$。
\end{proof}

当 $P$ 超过 $D_0$ 时，如例题中的镧（其中 $D_0 =
0.0074$，$P = 0.025$），残余体比其他矿物更快地失去最佳寄主矿物单斜辉石，
此时熔体比实比熔融的预测略为富集，
对镧而言，在 $F = 0.1$ 时偏差不到 \qty{2}{\percent}（练习~6），但
残余体在 $F = 0.12/0.45 \approx 0.27$ 时耗尽单斜辉石，超过此值后，
矿物与分配系数都会发生变化。
% ledger: white13:ex73

\section{分离熔融与动态熔融}\label{sec:b3:trace-element-modelling:fractional}

橄榄岩中的熔体沿颗粒边缘聚集，
一旦相互连通即可逸出。与\omterm{thm:b3:trace-element-modelling:batch}{批式熔融}相对的端元
则在熔体一经形成时即将其全部移出。

\begin{theorem}[分离熔融]\label{thm:b3:trace-element-modelling:fractional}
在\emph{分离熔融}\index{分离熔融}中，每一增量的熔体
仅在形成的瞬间与残余体处于平衡，随即离开。
当 $D$ 为常数时，熔体总质量 $F$ 被抽出后，
残余体与此时新形成的熔体满足
\[
\frac{C_S}{C_0} = (1 - F)^{1/D - 1}, \qquad
\frac{C_L}{C_0} = \frac{1}{D}(1 - F)^{1/D - 1}.
\]
\end{theorem}

\begin{proof}
残余体质量为 $1 - F$，浓度为 $C_S$。当质量为
$\mathrm dF$ 的熔体以 $C_L = C_S/D$ 的浓度离开时，残余体中的该元素
减少 $\mathrm d[(1 - F)C_S] = -(C_S/D)\,\mathrm dF$。展开得
$(1 - F)\,\mathrm dC_S = C_S(1 - 1/D)\,\mathrm dF$，故
$\mathrm dC_S/C_S = (1/D - 1)\,\mathrm d[-\ln(1 - F)]$。积分，由
$C_S = C_0$（$F = 0$ 时）出发，得 $\ln(C_S/C_0) = (1/D - 1)\ln(1 - F)$，并
$C_L = C_S/D$。
\end{proof}

对于不相容元素，$1/D - 1$ 很大且为正，残余体
在熔融量仅为 $D$ 的若干倍时即被耗尽，瞬时熔体
由 $C_0/D$ 降至几乎为零。\omterm{thm:b3:trace-element-modelling:fractional}{分离熔融}的残余体远比
\omterm{thm:b3:trace-element-modelling:batch}{批式熔融}的残余体贫化。这些熔体并非逐一喷出，
而是汇集于河道并聚集成池，下文定义的\omterm{def:b3:trace-element-modelling:fractional}{聚集分离熔体}
在图中以虚线画出。

\begin{omfigure}
\begin{tikzpicture}
\begin{groupplot}[group style={group size=2 by 1, horizontal sep=1.6cm},
  omrecord, width=0.47\linewidth, height=0.44\linewidth, title style={font=\small},
  xlabel={熔体分数 $F$}, xmin=0, xmax=0.3, ymode=log, ymin=0.01, ymax=150,
  xtick={0,0.1,0.2,0.3}, xticklabels={0,0.1,0.2,0.3},
  ytick={0.01,0.1,1,10,100}, yticklabels={0.01,0.1,1,10,100},
  legend style={font=\scriptsize, legend cell align=left}]
\nextgroupplot[title={(a) 批式（实线），聚集分离（虚线）},
  ylabel={$C_L/C_0$}, legend style={at={(0.97,0.97)}, anchor=north east}]
% computed in figdata/bachelor-3/melting_models.py
\addplot[omDef, thick] table[x=f, y=b001] {figdata/out/bachelor-3-melting_models-curves.dat};
\addlegendentry{$D = 0.01$}
\addplot[red!70!black, thick] table[x=f, y=b01] {figdata/out/bachelor-3-melting_models-curves.dat};
\addlegendentry{$D = 0.1$}
\addplot[black, thick] table[x=f, y=b1] {figdata/out/bachelor-3-melting_models-curves.dat};
\addlegendentry{$D = 1$}
\addplot[omLithLime!70!black, thick] table[x=f, y=b10] {figdata/out/bachelor-3-melting_models-curves.dat};
\addlegendentry{$D = 10$}
\addplot[omDef, thick, densely dashed] table[x=f, y=a001] {figdata/out/bachelor-3-melting_models-curves.dat};
\addplot[red!70!black, thick, densely dashed] table[x=f, y=a01] {figdata/out/bachelor-3-melting_models-curves.dat};
\addplot[omLithLime!70!black, thick, densely dashed] table[x=f, y=a10] {figdata/out/bachelor-3-melting_models-curves.dat};
\nextgroupplot[title={(b) 分离熔融，瞬时}]
\addplot[omDef, thick] table[x=f, y=i001] {figdata/out/bachelor-3-melting_models-curves.dat};
\addplot[red!70!black, thick] table[x=f, y=i01] {figdata/out/bachelor-3-melting_models-curves.dat};
\addplot[black, thick] table[x=f, y=i1] {figdata/out/bachelor-3-melting_models-curves.dat};
\addplot[omLithLime!70!black, thick] table[x=f, y=i10] {figdata/out/bachelor-3-melting_models-curves.dat};
\end{groupplot}
\end{tikzpicture}

\par\smallskip{\small 微量元素在熔体中的富集度随熔体分数
的变化，对应四个总分配系数。(a) 批式熔体（实线）
及\omterm{def:b3:trace-element-modelling:fractional}{聚集分离熔体}（虚线；$D = 1$ 时两者均保持为 1）。(b) \omterm{thm:b3:trace-element-modelling:fractional}{分离熔融}的
瞬时熔体：$D = 0.01$ 的元素在
熔融量达百分之几后即从残余体中耗尽。}\label{fig:b3:trace-element-modelling:curves}
\end{omfigure}

\begin{definition}[聚集分离熔体]\label{def:b3:trace-element-modelling:fractional}
称\emph{聚集分离熔体}\index{聚集分离熔体}为
\omterm{thm:b3:trace-element-modelling:fractional}{分离熔融}中从 $0$ 到
$F$ 阶段所有瞬时熔体的混合物，汇集后不再与残余体发生反应。
\end{definition}

\begin{proposition}[聚集熔体的成分]
\label{prop:b3:trace-element-modelling:aggregate}
设 $D$ 为常数，则\omterm{def:b3:trace-element-modelling:fractional}{聚集分离熔体}满足
\[
\frac{\bar C_L}{C_0} = \frac{1 - (1 - F)^{1/D}}{F}.
\]
\end{proposition}

\begin{proof}
被提取的元素等于源区减去残余体，
$C_0 - (1 - F)C_S = C_0[1 - (1 - F)^{1/D}]$，由\cref{thm:b3:trace-element-modelling:fractional}可得，
它均匀分布在质量为 $F$ 的熔体中。
\end{proof}

当 $(1 - F)^{1/D}$ 很小时，聚集熔体几乎包含了该元素的全部，
于是 $\bar C_L \approx C_0/F$，与 $D$ 无关：对于极不相容元素，
\omterm{def:b3:trace-element-modelling:fractional}{聚集分离熔体}与小 $D$ 的\omterm{thm:b3:trace-element-modelling:batch}{批式熔融}结果一致，两者
都度量 $F$。

实际熔融介于两个端元之间。残余体会在
其孔隙中保留一部分熔体，离开的熔体则是超过该孔隙度的那部分。

\begin{definition}[动态熔融]\label{def:b3:trace-element-modelling:dynamic}
在\emph{动态熔融}\index{动态熔融}（又称连续熔融）中，
残余体的孔隙中保留恒定质量分数 $\phi$ 的熔体，
与其晶体处于平衡，超出该量产生的熔体
全部被提取。
\end{definition}

\begin{proposition}[动态熔融]\label{prop:b3:trace-element-modelling:dynamic}
孔隙度为 $\phi$ 的\omterm{def:b3:trace-element-modelling:dynamic}{动态熔融}即为\omterm{thm:b3:trace-element-modelling:fractional}{分离熔融}，其有效
分配系数为 $D' = (1 - \phi)D + \phi$，其中 $F$ 为
被提取熔体的质量分数：
\[
\frac{C_L}{C_0} = \frac{1}{D'}(1 - F)^{1/D' - 1}.
\]
\end{proposition}

\begin{proof}
残余体由占比 $1 - \phi$ 的晶体与
占比 $\phi$ 的孔隙熔体组成，二者均与 $C_L$ 处的熔体平衡，因此
$C_S = [(1 - \phi)D + \phi]C_L = D'C_L$。于是可直接应用
\cref{thm:b3:trace-element-modelling:fractional}的证明，取 $D'$。
\end{proof}

当 $\phi = 0$ 时，该模型为\omterm{thm:b3:trace-element-modelling:fractional}{分离熔融}；当 $\phi \to 1$ 时，
它趋近于\omterm{thm:b3:trace-element-modelling:batch}{批式熔融}。孔隙度为百分之一时，第一批熔体的富集程度
被限制在 $1/\phi = 100$，无论 $D$ 多小：
孔隙熔体起到了缓冲残余体的作用。

\begin{omfigure}
\begin{tikzpicture}[font=\small, >=Stealth]
\tikzset{rock/.style={draw, fill=omLithLime!25, minimum width=2.3cm, minimum height=2.0cm},
  melt/.style={draw, fill=red!25}}
% (a) batch
\node[rock] (a) at (0,0) {};
\foreach \x/\y in {-0.7/0.5, 0.3/0.6, -0.2/-0.1, 0.7/-0.4, -0.6/-0.6, 0.4/0.0}
  \fill[red!55] (\x,\y) circle (0.09);
\node[align=center, font=\scriptsize] at (0,-1.45) {残余体 $1-F$，含 $DC_L$\\熔体 $F$，含 $C_L$};
\draw[->, thick, red!70!black] (0,1.05) -- (0,1.7) node[above, font=\scriptsize] {一次性全部离开};
\node[font=\small] at (0,-2.25) {(a) 批式};
% (b) fractional
\node[rock] (b) at (4,0) {};
\draw[->, thick, red!70!black] (4,1.05) -- (4,1.7);
\foreach \y/\s in {1.85/0.06, 2.05/0.08, 2.27/0.1}
  \fill[red!55] (4,\y) circle (\s);
\node[align=center, font=\scriptsize] at (4,-1.45) {每一小份形成后 \\立即离开};
\node[font=\small] at (4,-2.25) {(b) 分离};
% (c) dynamic
\node[rock] (c) at (8,0) {};
\foreach \x/\y in {7.4/0.5, 8.4/0.3, 7.8/-0.5, 8.6/-0.6}
  \fill[red!55] (\x,\y) circle (0.06);
\draw[->, thick, red!70!black] (8,1.05) -- (8,1.7);
\foreach \y/\s in {1.85/0.06, 2.05/0.08, 2.27/0.1}
  \fill[red!55] (8,\y) circle (\s);
\node[align=center, font=\scriptsize] at (8,-1.45) {孔隙保留 $\phi$；\\超出部分离开};
\node[font=\small] at (8,-2.25) {(c) 动态};
\end{tikzpicture}

\par\smallskip{\small 熔体离开熔融源区的三种方式。(a)
批式：熔体与整个残余体保持接触，之后整体离开。(b) 分离：
每个增量立即离开；汇集的各增量形成
\omterm{def:b3:trace-element-modelling:fractional}{聚集分离熔体}。(c) 动态：残余体的孔隙中保留
小质量分数 $\phi$ 的熔体。}\label{fig:b3:trace-element-modelling:modes}
\end{omfigure}

\section{熔融之后的结晶}\label{sec:b3:trace-element-modelling:crystallisation}

离开源区的熔体很少原样喷出。它在上升途中会先结晶出
橄榄石；随后结晶出辉石和斜长石；它可能同化围岩，
也可能与其他批次的熔体混合。本科二年级分册中的规律依次适用：
熔岩是源区中占比 $F$ 的熔体，随后成为一种液体，
该液体经分离结晶保留了其质量的 $f$，于是
对于在熔融残余体中的分配系数为 $D$、在结晶矿物中为 $D_c$ 的元素，
则有
\[
C_{\text{熔岩}} = C_0\,\frac{1}{D + F(1 - D)}\,f^{\,D_c - 1}.
\]
第二个因子对结晶矿物所排斥的每种元素都相同，
即 $D_c \approx 0$：它使所有这类元素都按 $1/f$ 变化。这类元素之间的比值
因此由源区决定，一种已失去橄榄石的熔岩，
其 La/Yb 保持原始熔体的值，而镧的含量则升高。由浓度反演 $F$
还需要 $f$；由比值反演源区则
不需要。斜长石是稀土元素中的例外，它会富集铕，
其去除则留下本科二年级分册中的铕异常。

\section{解读稀土元素配分模式}\label{sec:b3:trace-element-modelling:patterns}

尖晶石二辉橄榄岩与石榴子石二辉橄榄岩的熔体共同拥有
轻稀土元素，两种残余体都不保留它们；而二者的重稀土元素不同，
后者由石榴子石保留。在球粒陨石标准化图解上，石榴子石熔体
的曲线较陡，重稀土端呈平底；尖晶石熔体的曲线则较平缓。
随着 $F$ 减小，曲线的陡度增大。

\begin{omfigure}
\begin{tikzpicture}
\begin{axis}[omrecord, width=0.8\linewidth, height=0.5\linewidth, ymode=log,
  xmin=56.5, xmax=71.5, ymin=1, ymax=1000, xtick={57,58,60,62,63,64,65,70,71},
  xticklabels={La,Ce,Nd,Sm,Eu,Gd,Tb,Yb,Lu}, ytick={1,10,100,1000}, yticklabels={1,10,100,1000},
  ylabel={样品 / CI 球粒陨石},
  legend style={font=\scriptsize, at={(0.98,0.98)}, anchor=north east,
  legend cell align=left, legend columns=2}]
% computed in figdata/bachelor-3/melting_models.py
\addplot[omDef, thick, mark=*, mark size=1.3pt] table[x=z, y=sp01] {figdata/out/bachelor-3-melting_models-ree.dat};
\addlegendentry{尖晶石，$F = 1\,\%$}
\addplot[red!70!black, thick, mark=square*, mark size=1.3pt] table[x=z, y=gt01] {figdata/out/bachelor-3-melting_models-ree.dat};
\addlegendentry{石榴子石，$F = 1\,\%$}
\addplot[omDef, thick, densely dashed, mark=*, mark size=1.3pt] table[x=z, y=sp05] {figdata/out/bachelor-3-melting_models-ree.dat};
\addlegendentry{尖晶石，$5\,\%$}
\addplot[red!70!black, thick, densely dashed, mark=square*, mark size=1.3pt] table[x=z, y=gt05] {figdata/out/bachelor-3-melting_models-ree.dat};
\addlegendentry{石榴子石，$5\,\%$}
\addplot[omDef, thick, densely dotted, mark=*, mark size=1.3pt] table[x=z, y=sp20] {figdata/out/bachelor-3-melting_models-ree.dat};
\addlegendentry{尖晶石，$20\,\%$}
\addplot[red!70!black, thick, densely dotted, mark=square*, mark size=1.3pt] table[x=z, y=gt20] {figdata/out/bachelor-3-melting_models-ree.dat};
\addlegendentry{石榴子石，$20\,\%$}
\addplot[black, thin, mark=none] table[x=z, y=pm] {figdata/out/bachelor-3-melting_models-ree.dat};
\addlegendentry{源区}
\end{axis}
\end{tikzpicture}

\par\smallskip{\small \omterm{thm:b3:trace-element-modelling:batch}{批式熔融}熔体的球粒陨石标准化稀土元素配分模式，
原始地幔源区，分别来自尖晶石二辉橄榄岩（蓝色圆点）与
石榴子石二辉橄榄岩（红色方块），取三种熔体分数；源区为
细线。仅有列表分配系数的九种稀土元素被
绘出。}\label{fig:b3:trace-element-modelling:ree}
\end{omfigure}
% ledger: budahn91:d-ree, ms95:pm-trace, lodders:ci-table, white13:ex73, rfr25:gkr-5gpa

地幔中的富铝矿物在约 60
至 \qty{90}{km} 深度处由石榴子石转为尖晶石。重稀土平坦且含量低、轻重稀土斜率陡峭的熔体，
因此其残余体中留有石榴子石，
熔融至少始于该深度。
% ledger: white13:sp-gt

一个配分模式同时混有两个未知量，即 $F$ 与源区。
将两种元素相对作图，即可把二者分开。

\begin{definition}[比值–比值图]\label{def:b3:trace-element-modelling:ratios}
称\emph{比值–比值图}\index{比值–比值图}为这样一种图：对一组熔岩，
以高度不相容元素 $H$ 与较不相容元素 $M$ 的比值，
对 $H$ 的浓度（或另一比值）作图。
部分熔融、结晶与混合在其上
画出形状各异的线。
\end{definition}

\begin{proposition}[批式熔融线]\label{prop:b3:trace-element-modelling:line}
对于同一源区的\omterm{thm:b3:trace-element-modelling:batch}{批式熔融}熔体，若 $D_H = 0$ 且 $D_M$ 为常数，则点
$(C_H, C_H/C_M)$ 位于直线
\[
\frac{C_H}{C_M} = \frac{D_M}{C_0^M}\,C_H + \frac{C_0^H(1 - D_M)}{C_0^M}.
\]
\end{proposition}

\begin{proof}
由 $D_H = 0$ 得 $C_H = C_0^H/F$，故 $1/F = C_H/C_0^H$。于是
$C_H/C_M = (C_0^H/F)\,(D_M + F(1 - D_M))/C_0^M = C_0^HD_M/(FC_0^M) +
C_0^H(1 - D_M)/C_0^M$，而第一项为 $(D_M/C_0^M)\,C_H$。
\end{proof}

分离结晶中，若矿物排斥这两种元素，熔岩
在此图上水平移动，因比值不变；熔融则使其
沿斜线移动，斜线斜率随 $D_M$ 增大；两种熔体混合则画出
双曲线。含石榴子石的源区具有较大的 $D_{\mathrm{Yb}}$，因而
在 La/Yb 对 La 图上的熔融线较陡。

\begin{method}[反演两个熔岩]\label{met:b3:trace-element-modelling:invert}
（1）检验这些熔岩是否同源，且 $H$ 远比 $M$ 不相容；
（2）对每个熔岩作结晶校正，或只用比值。
（3）利用两个或更多熔岩，拟合直线 $C_H/C_M = aC_H + b$。
（4）则 $a = D_M/C_0^M$，$b = C_0^H(1 - D_M)/C_0^M$。若 $C_0^M$ 已设定，则 $D_M =
aC_0^M$，$C_0^H = bC_0^M/(1 - D_M)$。（5）各熔岩的熔体分数为
$F = C_0^H/C_H$。
\end{method}

\begin{omfigure}
\begin{tikzpicture}
\begin{axis}[omrecord, width=0.78\linewidth, height=0.52\linewidth,
  xmode=log, ymode=log, xmin=2, xmax=60, ymin=2, ymax=100,
  xlabel={La (\unit{ppm})}, ylabel={La/Yb},
  xtick={2,5,10,20,50}, xticklabels={2,5,10,20,50},
  ytick={2,5,10,20,50,100}, yticklabels={2,5,10,20,50,100},
  legend style={font=\scriptsize, at={(0.03,0.97)}, anchor=north west,
  legend cell align=left}]
% computed in figdata/bachelor-3/melting_models.py
\addplot[omDef, thick] table[x=la_sp, y=lay_sp] {figdata/out/bachelor-3-melting_models-traj.dat};
\addlegendentry{尖晶石二辉橄榄岩}
\addplot[red!70!black, thick] table[x=la_gt, y=lay_gt] {figdata/out/bachelor-3-melting_models-traj.dat};
\addlegendentry{石榴子石二辉橄榄岩}
\addplot[omDef, only marks, mark=|, mark size=3pt, forget plot] table[x=la_sp, y=lay_sp] {figdata/out/bachelor-3-melting_models-ticks.dat};
\addplot[red!70!black, only marks, mark=|, mark size=3pt, forget plot] table[x=la_gt, y=lay_gt] {figdata/out/bachelor-3-melting_models-ticks.dat};
\addplot[black, only marks, mark=diamond*, mark size=2.8pt] table[x=la, y=layb] {figdata/out/bachelor-3-melting_models-lavas.dat};
\addlegendentry{熔岩 A 与 B}
\node[font=\scriptsize, anchor=north] at (axis cs:9.6,5.2) {A};
\node[font=\scriptsize, anchor=south east] at (axis cs:15.2,36) {B};
\node[font=\scriptsize, omDef, anchor=north west] at (axis cs:36.5,17.5) {1\,\%};
\node[font=\scriptsize, omDef, anchor=north west] at (axis cs:11.3,6.2) {5\,\%};
\node[font=\scriptsize, omDef, anchor=north west] at (axis cs:3.14,2.55) {20\,\%};
\node[font=\scriptsize, red!70!black, anchor=south east] at (axis cs:29.5,66) {1\,\%};
\node[font=\scriptsize, red!70!black, anchor=south east] at (axis cs:10.5,24) {5\,\%};
\node[font=\scriptsize, red!70!black, anchor=south east] at (axis cs:3.09,7.0) {20\,\%};
\end{axis}
\end{tikzpicture}

\par\smallskip{\small 原始地幔源区的\omterm{thm:b3:trace-element-modelling:batch}{批式熔融}熔体在 La/Yb
对 La 图上的变化，分别针对尖晶石二辉橄榄岩与石榴子石二辉橄榄岩，刻度位于 $F = 1$、
2、5、10 与 \qty{20}{\percent}，并标出周末问题中的两个虚构熔岩。
每个熔岩仅位于一条曲线上。}\label{fig:b3:trace-element-modelling:ratios}
\end{omfigure}

在对数坐标上，命题所述的直线会发生弯曲，并且当
$D_{\mathrm{La}} \ne 0$ 时，它在线性坐标上只是近似为直线；图中的曲线是精确的。
尖晶石源区所产生的熔体，其 La/Yb 不可能高于
$(C_0^{\mathrm{La}}/D_{\mathrm{La}})/(C_0^{\mathrm{Yb}}/D_{\mathrm{Yb}})$,
对原始地幔约为 40；石榴子石源区可达约 120。

\section{练习}

\begin{exercise}[$\star$]\label{exo:b3:trace-element-modelling:1}
已知钐的矿物分配系数：单斜辉石中为 0.40，石榴子石中为 0.11，
在橄榄石、斜方辉石和尖晶石中为 0.016。计算钐在
本章尖晶石二辉橄榄岩（\qty{12}{\percent} 单斜辉石）
与石榴子石二辉橄榄岩（\qty{19.5}{\percent} 单斜辉石、
\qty{16.7}{\percent} 石榴子石）中的总分配系数。
\end{exercise}
% ledger: budahn91:d-ree

\begin{exercise}[$\star$]\label{exo:b3:trace-element-modelling:2}
某元素 $D = 0.05$，其在源区中的含量为 \qty{1}{ppm}。计算
其在 $F = 0.02$ 与 $0.10$ 时\omterm{thm:b3:trace-element-modelling:batch}{批式熔融}熔体中的浓度。
\end{exercise}

\begin{exercise}[$\star$]\label{exo:b3:trace-element-modelling:3}
\omterm{thm:b3:trace-element-modelling:batch}{批式熔融}能使某元素达到的最大富集比 $C_L/C_0$ 为多少？
对 $D = 0.01$ 的元素而言；对 $D = 0$ 的元素，在
$F = \qty{5}{\percent}$ 时的最大富集比又为多少？
\end{exercise}

\begin{exercise}[$\star$]\label{exo:b3:trace-element-modelling:4}
某熔岩中 La 为 \qty{30}{ppm}，Yb 为 \qty{2.5}{ppm}。已知 CI 球粒陨石值
La 为 0.232、Yb 为 \qty{0.163}{ppm}，计算其 La/Yb 及其
球粒陨石标准化值 $(\mathrm{La}/\mathrm{Yb})_N$。
\end{exercise}
% ledger: lodders:ci-table

\begin{exercise}[$\star\star$]\label{exo:b3:trace-element-modelling:5}
对 $D = 0.01$ 的元素，比较经 \qty{5}{\percent} 的
\omterm{thm:b3:trace-element-modelling:batch}{批式熔融}与\omterm{thm:b3:trace-element-modelling:fractional}{分离熔融}后的残余
成分。此时瞬时分离熔体中该元素的含量为多少？
\end{exercise}

\begin{exercise}[$\star\star$]\label{exo:b3:trace-element-modelling:6}
在本章的例题中（模式含量分别为 58、27、12、\qty{3}{\percent}；
熔融比例分别为 20、25、45、\qty{10}{\percent}，对应橄榄石、
斜方辉石、单斜辉石与尖晶石），镧的分配系数依次为
0.0003、0.002、0.053 与 0.01。计算 $D_0$ 与 $P$，再求
镧在 $F = 0.02$ 与 $0.10$ 时批式熔体中的富集度，分别按非实比与实比熔融计算。
\end{exercise}
% ledger: white13:ex73

\begin{exercise}[$\star\star$]\label{exo:b3:trace-element-modelling:7}
在 $F = 0.10$ 时，比较批式熔体与\omterm{def:b3:trace-element-modelling:fractional}{聚集分离熔体}，
分别针对 $D = 0.1$ 与 $D = 0.01$。哪一者差异更大？为什么？
\end{exercise}

\begin{exercise}[$\star\star$]\label{exo:b3:trace-element-modelling:8}
某元素 $D = 0.001$，以 $\phi = 0.01$ 进行\omterm{def:b3:trace-element-modelling:dynamic}{动态熔融}。计算
$D'$、最先抽出的熔体的富集度，以及 $F = 0.01$ 时抽出的熔体的富集度；
并与\omterm{thm:b3:trace-element-modelling:fractional}{分离熔融}进行比较。
\end{exercise}

\begin{exercise}[$\star\star$]\label{exo:b3:trace-element-modelling:9}
在本章的尖晶石二辉橄榄岩中，$D_{\mathrm{La}} = 0.0078$ 与
$D_{\mathrm{Sm}} = 0.062$。\omterm{thm:b3:trace-element-modelling:batch}{批式熔融}使 La/Sm 比源区值
分别在 $F = 0.01$ 与 $F = 0.10$ 时提高多少倍？这说明
用比值估计 $F$ 有何局限？
\end{exercise}

\begin{exercise}[$\star\star\star$]\label{exo:b3:trace-element-modelling:10}
证明\omterm{def:b3:trace-element-modelling:fractional}{聚集分离熔体}是从 0 到 $F$ 的瞬时熔体的
平均值，并证明它趋于 $C_0/F$ 的条件为 $(1 - F)^{1/D} \to 0$。
对 $D = 0.01$，从何种 $F$ 开始它与 $C_0/F$ 之差在 \qty{1}{\percent} 以内？
\end{exercise}

\begin{exercise}[$\star\star\star$]\label{exo:b3:trace-element-modelling:11}
两个同源熔岩 P 与 Q 的 La 分别为 20 与 \qty{10}{ppm}，La/Yb 分别为 10 与 6。
假设镧完全不相容，且源区含有
\qty{0.44}{ppm} Yb，应用本章方法求 $D_{\mathrm{Yb}}$、
源区中 La 的含量以及每个熔岩的熔体分数。
\end{exercise}

\begin{exercise}[$\star\star\star$]\label{exo:b3:trace-element-modelling:12}
在练习~6 的\omterm{def:b3:trace-element-modelling:nonmodal}{非实比熔融}中，各矿物在何种熔体分数下
耗尽？计算镧在 $F = 0.2$ 时残余物的
总分配系数，并说明为何本章公式在
首次耗尽之后不能使用。
\end{exercise}

\section{问题：阿尔德伦岛的两种熔岩}

\begin{problem}[{周末问题——两种虚构玄武岩的源区、深度与熔融程度
由其镧与镱含量推定}]\label{pb:b3:trace-element-modelling:1}
在虚构的火山岛阿尔德伦岛上，熔岩 A 含 La \qty{9.6}{ppm}
与 Yb \qty{1.7}{ppm}，熔岩 B 含 La \qty{15.6}{ppm} 与 Yb \qty{0.45}{ppm}。
两者均取为原始地幔源区的原生熔体，该源区中
La 为 \qty{0.648}{ppm}、Yb 为 \qty{0.441}{ppm}（CI 球粒陨石：La 0.232，Yb
\qty{0.163}{ppm}）。矿物分配系数：La 在单斜辉石中为 0.058，在
石榴子石、橄榄石、斜方辉石与尖晶石中均为 0.001；Yb 在单斜辉石中为 0.58，
在石榴子石中为 4.5，在橄榄石、斜方辉石与尖晶石中为 0.16。
尖晶石二辉橄榄岩由 \qty{58}{\percent} 橄榄石、\qty{27}{\percent} 斜方辉石、\qty{12}{\percent}
单斜辉石与 \qty{3}{\percent} 尖晶石组成，石榴子石二辉橄榄岩由 \qty{54.1}{\percent} 橄榄石、\qty{9.8}{\percent}
斜方辉石、\qty{19.5}{\percent} 单斜辉石与 \qty{16.7}{\percent} 石榴子石组成。本题中熔融为实比熔融。
% ledger: ms95:pm-trace, lodders:ci-table, budahn91:d-ree, white13:ex73, rfr25:gkr-5gpa

\medskip\noindent\textbf{第一部分——两种源区。}
\begin{enumerate}
\item 计算尖晶石二辉橄榄岩的 $D_{\mathrm{La}}$ 与 $D_{\mathrm{Yb}}$。
\item 计算石榴子石二辉橄榄岩的上述两个系数。为便于后续计算，将四个值均
保留两位有效数字，用于本题其余部分。
\item 两种源区之间，哪一个分配系数差异最大？原因何在？
镱在石榴子石二辉橄榄岩中是否为相容元素？
\item 计算源区的球粒陨石标准化
值 $(\mathrm{La}/\mathrm{Yb})_N$，其标准化图解呈何种形态？
\end{enumerate}

\medskip\noindent\textbf{第二部分——熔岩 A 的\omterm{thm:b3:trace-element-modelling:batch}{批式熔融}。}
\begin{enumerate}[resume]
\item 计算尖晶石二辉橄榄岩在 $F = 1$、5 和
\qty{10}{\percent} 时批式熔体中的 La。
\item 求使熔岩 A 具有其镧含量的熔体分数。
\item 预测其 Yb 及 La/Yb，并与熔岩 A 比较。
\item 计算残留体中的 La 和 Yb，以及源区中
留存于残留体中的镧所占的比例。
\item 若将镧视为完全不相容，
所得的熔体分数是多少？
\end{enumerate}

\medskip\noindent\textbf{第三部分——熔岩 A 的\omterm{thm:b3:trace-element-modelling:fractional}{分离熔融}。}
\begin{enumerate}[resume]
\item 计算 $F = 0.1$、1 和
\qty{6}{\percent} 时瞬时分离熔体中的 La。
\item 计算 $F = \qty{6.0}{\percent}$ 时\omterm{def:b3:trace-element-modelling:fractional}{聚集分离熔体}中的 La。
\item 求使聚集熔体具有熔岩 A 镧含量的熔体分数。
为何它与第 9 题的答案如此接近？
\item 将残留体的镧含量与第 8 题的结果进行比较。
\item 你保留哪一个 $F$ 的估计值？
\end{enumerate}

\medskip\noindent\textbf{第四部分——熔岩 B 与石榴子石。}
\begin{enumerate}[resume]
\item 计算两种熔岩的 La/Yb 及 $(\mathrm{La}/\mathrm{Yb})_N$。
\item 尖晶石二辉橄榄岩的批式熔体
所能达到的最大 La/Yb 是多少？
\item 在何种熔体分数下，尖晶石二辉橄榄岩熔体的
La/Yb 与熔岩 B 相同？其中含镧多少？请作出结论。
\item 对石榴子石二辉橄榄岩，由熔岩 B 的镧含量
求其熔体分数，并预测其 Yb。
\item 两种源区分别说明 A 与 B 开始熔融的深度
如何？熔岩 A 能否来自石榴子石二辉橄榄岩？
\item 对于石榴子石源区，计算\omterm{thm:b3:trace-element-modelling:batch}{批式熔融}线 La/Yb $= a\,
\mathrm{La} + b$ 的斜率与截距，取 $D_{\mathrm{La}}$ 为零，
并用熔岩 B 进行检验。
\item 若熔岩 A 在喷发前失去了 \qty{10}{\percent} 的质量作为橄榄石，
其 La/Yb 与熔体分数将如何变化？
\item 回答本章开篇的问题，针对这两种熔岩作答。
\item 求熔岩 A 的熔融程度。
\end{enumerate}
\end{problem}
